\chapter{Deadlocks}
\label{ch:deadlocks}

Locks fix races and create the next problem: threads that wait for each
other forever. With two teams adding locks to shared SWEB structures ---
the thread list, the address space, the scheduler --- your group has to
agree on how locks are taken. This chapter gives the theory (Coffman's
conditions, graphs, prevention, avoidance, detection) and the practice
(lock ordering, validators), plus the relatives of deadlock: livelock,
starvation and priority inversion.

\section{Definition and the four conditions}

A set of threads is \emph{deadlocked} if every thread in the set waits
for an event that only another thread in the set can cause. The smallest
example: thread 1 holds lock A and waits for B, thread 2 holds B and waits
for A. A deadlock is possible only if all four \emph{Coffman conditions}
hold:

\begin{enumerate}
  \item \textbf{Mutual exclusion}: a resource is held by at most one
    thread.
  \item \textbf{Hold and wait}: a thread holds resources while waiting for
    more.
  \item \textbf{No preemption}: resources cannot be taken away; the holder
    releases them voluntarily.
  \item \textbf{Circular wait}: a cycle T$_1$ $\to$ T$_2$ $\to$ \dots $\to$
    T$_1$, each waiting for a resource the next one holds.
\end{enumerate}

Break any one and deadlock becomes impossible.

\section{Resource-allocation and wait-for graphs}

Draw threads as circles and resources as squares; an edge resource
$\to$ thread means ``held by'', thread $\to$ resource means ``waits
for''. With one instance per resource, a \emph{cycle means deadlock}.
Collapsing the resources gives the \emph{wait-for graph}: T$_i$ $\to$
T$_j$ if T$_i$ waits for something T$_j$ holds.

\begin{figure}[htbp]
\centering
\begin{tikzpicture}[font=\small,node distance=1.4cm,
    t/.style={circle,draw,minimum size=0.8cm},
    r/.style={rectangle,draw,minimum size=0.7cm,fill=black!6},
    >=Latex]
  \node[t] (t1) {T$_1$};
  \node[r,right=of t1] (b) {B};
  \node[t,right=of b] (t2) {T$_2$};
  \node[r,below=1.1cm of t2] (c) {C};
  \node[t,left=of c] (t3) {T$_3$};
  \node[r,below=1.1cm of t1] (a) {A};
  \draw[->] (a) -- node[left] {held} (t1);
  \draw[->,dashed] (t1) -- node[above] {waits} (b);
  \draw[->] (b) -- node[above] {held} (t2);
  \draw[->,dashed] (t2) -- node[right] {waits} (c);
  \draw[->] (c) -- node[below] {held} (t3);
  \draw[->,dashed,warnrule,thick] (t3) -- node[below] {waits?} (a);
\end{tikzpicture}
\caption{T$_1$ holds A and waits for B, T$_2$ holds B and waits for C, T$_3$
holds C. No cycle: T$_3$ can finish. If T$_3$ now requests A, the cycle
T$_1\to$T$_2\to$T$_3\to$T$_1$ closes: deadlock.}
\label{fig:rag}
\end{figure}

\section{Strategies}

\paragraph{Prevention} --- make one condition impossible:
\begin{itemize}
  \item mutual exclusion: share instead (read-only data, lock-free
    structures) --- only for some resources;
  \item hold and wait: acquire everything at once, or release everything
    before waiting (trylock and back off --- risks livelock);
  \item no preemption: take resources away (roll back a transaction) ---
    rarely possible for locks;
  \item circular wait: \textbf{a global lock order}. If every thread takes
    locks in the same order, no cycle can form. This is the standard
    technique in kernels.
\end{itemize}

\paragraph{Avoidance} --- grant a request only if the system stays in a
\emph{safe state}: one from which some order lets every thread finish,
even if all of them ask for their declared maximum. The \emph{banker's
algorithm} checks this. Example with two resource types A (10 units) and
B (5 units):

\begin{center}
\begin{tabular}{lccc}
\toprule
 & \textbf{Allocation} & \textbf{Max} & \textbf{Need = Max -- Alloc} \\
\midrule
P$_0$ & (2, 1) & (7, 3) & (5, 2) \\
P$_1$ & (3, 1) & (4, 2) & (1, 1) \\
P$_2$ & (2, 2) & (5, 4) & (3, 2) \\
\bottomrule
\end{tabular}
\end{center}

Available = (10, 5) -- (7, 4) = (3, 1). Safety check: P$_1$'s need (1, 1)
$\le$ (3, 1): it can finish and return its allocation, Work = (6, 2);
then P$_2$ (3, 2) $\le$ (6, 2), Work = (8, 4); then P$_0$ (5, 2). Safe,
sequence P$_1$, P$_2$, P$_0$.
\begin{itemize}
  \item P$_0$ requests (1, 0): pretend to grant, Available = (2, 1), P$_0$'s
    need (4, 2). P$_1$, then P$_2$, then P$_0$ can still finish: safe,
    grant.
  \item P$_2$ requests (1, 1) instead: Available would be (2, 0). No
    thread's need fits into (2, 0) --- unsafe: P$_2$ must wait, although
    the units are free.
\end{itemize}
Avoidance needs to know maximum demands in advance; general-purpose
kernels therefore rarely use it, but it is a classic exam topic.

\paragraph{Detection and recovery} --- let deadlocks happen, find the
cycle in the wait-for graph, and break it (kill a thread, roll back).
Databases do this. A kernel can at least \emph{report} a deadlock instead
of hanging silently.

\paragraph{Ignore it} --- the ``ostrich algorithm''. For locks inside a
kernel, the real answer is prevention by lock ordering, backed by tools
that check the ordering.

\section{Lock ordering in practice}

\begin{itemize}
  \item Document a hierarchy: ``process lock before thread lock before
    address-space lock''. Code that holds a lower lock never acquires a
    higher one.
  \item Same-level locks (two accounts, two threads): order by a key ---
    the account number, the address of the lock object.
  \item Never call out into unknown code (callbacks, other subsystems)
    while holding a lock; you do not know which locks it takes.
  \item Keep critical sections small, and never sleep while holding a lock
    somebody else needs to wake you.
\end{itemize}

\paragraph{Validators.} Deadlocks depend on timing, tests rarely hit
them. A \emph{lock-order validator} records ``A was held while B was
acquired'' edges during a normal run and reports a cycle as soon as some
code path takes B then A --- even if the deadlock never happened. Linux
calls it \emph{lockdep}; ThreadSanitizer has one built in; \module{06},
part C builds one.

\section{Relatives of deadlock}

\begin{description}
  \item[Livelock] threads keep running but make no progress --- two
    threads that each trylock, fail, back off and retry in lockstep.
    Randomised back-off breaks the symmetry.
  \item[Starvation] the system makes progress, but one thread never gets
    its turn: an unfair lock, a reader-preferring rwlock, a low priority
    under strict priority scheduling (chapter~9).
  \item[Priority inversion] a high-priority thread waits for a lock held
    by a low-priority thread, while a medium-priority thread (which needs
    no lock) keeps the low one from running. The high thread effectively
    waits for the medium one. Mars Pathfinder (1997) reset itself because
    of this. \emph{Priority inheritance}: the lock holder temporarily
    runs with the priority of its highest waiter.
\end{description}

\begin{swebbox}
Every \code{Lock} knows its holder (\code{held\_by\_}); every
\code{Thread} keeps the locks it holds (\code{holding\_lock\_list\_}) and
the lock it waits for (\code{lock\_waiting\_on\_}). Before a thread goes
to sleep on a lock, \code{doChecksBeforeWaiting} can walk ``waits for lock
held by thread that waits for \dots'' and detect a real deadlock, and it
catches recursive locking. \code{\~{}Thread} asserts that a dying thread
holds no locks. Interrupt handlers must not take ordinary locks at all:
a handler waiting for a lock held by the thread it interrupted waits
forever --- a deadlock with a single thread.
\end{swebbox}

\begin{keybox}
\begin{itemize}
  \item Deadlock needs all four Coffman conditions; break one.
  \item Single-instance resources: a cycle in the wait-for graph is a
    deadlock.
  \item Kernels prevent deadlock with a documented lock order.
  \item Banker's algorithm: grant only if a safe sequence still exists.
  \item Validators find potential deadlocks from one run; SWEB's lock
    lists support similar checks.
  \item Livelock, starvation and priority inversion are the neighbours to
    recognise.
\end{itemize}
\end{keybox}
